% Part: first-order-logic
% Chapter: syntax-and-semantics
% Section: introduction

\documentclass[../../../include/open-logic-section]{subfiles}

\begin{document}

\olfileid{fol}{syn}{int}

\olsection{પરિચય}

પ્રથમ-ક્રમ તર્કશાસ્ત્રના સિદ્ધાંત અને
અધિસિદ્ધાંતનો વિકાસ કરવા માટે આપણે
પહેલાં તેની અભિવ્યક્તિઓની વાક્યરચના
અને અર્થવિચાર વ્યાખ્યાયિત કરવાં પડે.
પ્રથમ-ક્રમ તર્કશાસ્ત્રની અભિવ્યક્તિઓ
પદો અને !!{formula}s છે. પદો
!!{variable}s, !!{constant}s અને
!!{function}sમાંથી બને છે.
પછી !!^{formula}s, !!{predicate}s
સાથે પદો જોડીને રચાય છે
(તેનાથી સૌથી નાનાં, “આણ્વિક”
!!{formula}s બને છે); અને આણ્વિક
!!{formula}sમાંથી તાર્કિક સંયોજકો
અને પરિમાણકો વડે વધુ જટિલ
સૂત્રો રચી શકીએ છીએ.
રચના-નિયમો આપવાની ઘણી જુદી
રીતો છે; અહીં એક શક્ય
રીત આપીએ છીએ. બીજાં તંત્રો
જુદા સંકેતો પસંદ કરશે,
જુદા સંયોજક-ગણોને મૂળભૂત
ગણશે અને કૌંસનો જુદી
રીતે ઉપયોગ કરશે
(અથવા કહેવાતી પોલિશ
સંકેતપદ્ધતિની જેમ બિલકુલ
નહીં વાપરે). જોકે બધા
અભિગમોમાં સમાન વાત
એ છે કે રચના-નિયમો
પદો અને !!{formula}sના ગણને
\emph{અનુમાનાત્મક રીતે}
વ્યાખ્યાયિત કરે છે.
યોગ્ય રીતે કરાય તો
દરેક અભિવ્યક્તિ મૂળભૂત
રીતે માત્ર એક જ રીતે
બની શકે છે. અભિવ્યક્તિઓને
\emph{એકમાત્ર વાચનીય}
બનાવતી અનુમાનાત્મક
વ્યાખ્યાને લીધે એ જ
પદ્ધતિ વડે આ
અભિવ્યક્તિઓને અર્થ
આપી શકીએ છીએ.

અભિવ્યક્તિઓને અર્થ આપવાનું ક્ષેત્ર
અર્થવિચાર છે. અર્થવિચારનો કેન્દ્રસ્થ
ખ્યાલ !!a{structure}માં સંતોષનો છે.
!!^a{structure} ભાષાના પાયાના
ઘટકોને અર્થ આપે છે:
!!a{domain} એટલે વસ્તુઓનો
અરિક્ત ગણ. પરિમાણકોનું
અર્થઘટન આ વ્યાપકક્ષેત્ર
પર વ્યાપ્ત થતાં થાય છે;
!!{constant}sને તેમાંના ઘટકો,
!!{function}sને !!{domain}માંથી
પોતાનામાં જતાં વિધેયો અને
!!{predicate}sને !!{domain}
પરના સંબંધો નિયુક્ત થાય છે.
!!{domain} અને પાયાના
સંકેતભંડોળને કરેલી નિયુક્તિઓ
મળીને !!a{structure} બને છે.
!!^{variable}s, !!{formula}sમાં
આવી શકે છે; અર્થવિચાર
આપવા માટે તેમને પણ
!!{domain}ના !!{element}s
નિયુક્ત કરવા પડે છે—આ
ચલ-નિયુક્તિ છે. અંતે
સંતોષસંબંધ આ બધી બાબતોને
એકત્ર કરે છે. !!^a{formula},
!!a{structure}~$\Struct{M}$માં
ચલ-નિયુક્તિ~$s$ સાપેક્ષ
સંતોષાઈ શકે છે; તેને
$\Sat{M}{!A}[s]$ લખીએ છીએ.
આ સંબંધ પણ~$!A$ની
રચના પરના અનુમાનથી
વ્યાખ્યાયિત થાય છે:
સંયોજકોની સત્યતા-સારણીઓ
વાપરીને $!A \land !B$નો
સંતોષ, $!A$ અને~$!B$
સંતોષાય છે (કે નથી)
તેના આધારે વ્યાખ્યાયિત
થાય છે. પછી એવું
નીકળે છે કે
!!{formula}~$!A$
!!a{sentence} હોય,
એટલે તેમાં મુક્ત
ચલો ન હોય, તો
ચલ-નિયુક્તિનો ફેર
પડતો નથી. તેથી
!!{sentence}s કોઈ
!!{structure}sમાં
સાદી રીતે સંતોષાય
છે (કે નથી) એમ
કહી શકીએ છીએ.

વાક્યો માટેના સંતોષસંબંધ
$\Sat{M}{!A}$ના આધારે
પ્રામાણ્ય, ફલિતતા અને
સંતોષ્યતાના પાયાના
અર્થવિચારી ખ્યાલો
વ્યાખ્યાયિત કરી શકીએ.
દરેક સંરચના કોઈ
વાક્યને સંતોષે,
એટલે $\Entails !A$
હોય, તો તે પ્રમાણભૂત છે.
!!{sentence}sના ગણમાંથી
તે ફલિત થાય, એટલે
$\Gamma \Entails !A$,
જો $\Gamma$નાં
બધાં !!{sentence}sને
સંતોષતી દરેક
!!{structure},~$!A$ને
પણ સંતોષે. કોઈ એક
!!{structure} કોઈ
વાક્યગણમાંનાં બધાં
!!{sentence}sને
એકસાથે સંતોષે,
તો તે ગણ સંતોષ્ય છે.
!!{formula}s અનુમાનાત્મક
રીતે વ્યાખ્યાયિત છે
અને સંતોષ પણ
!!{formula}sની રચના
પરના અનુમાનથી
વ્યાખ્યાયિત થાય છે;
તેથી અર્થવિચારના
ગુણધર્મો સાબિત
કરવા અને આ
અર્થવિચારી ખ્યાલોને
પરસ્પર જોડવા આપણે
અનુમાન વાપરી શકીએ છીએ.




\end{document}
